\section{Polarizations and the common expansion}\label{sec:geometry}
Let $\rho$ denote projection off the curve on any stable-pair
parameter space, and $E$ the pushforward of the sheaf to the curve.
This projection on the support is finite, being affine and proper.
Purity and flatness of the family make $E$ a vector bundle of rank
$m$. We take the relative Euler characteristic on the proper torus-fixed
locus with its virtual structure sheaf, times the inverse moving
Euler $\lambda_{-1}((T^{\rm mov})^\vee)^{-1}$ and the ordered
determinant line of \cite{NO,W}. All obstruction theories are relative
to $B$. Equivariant $K$-theoretic virtual localization is used in the
sense of \cite{Q}. We give the two lines explicitly, suppressing
constant one-dimensional vector spaces with trivial character.

\begin{lemma}\label{lem:root}
Write $\mathcal D(E)=\det R\rho_*E$ and use angular brackets for the
Deligne pairing on the curve. With $N=tL$, the original polarization,
virtual tangent and modified determinant at Euler characteristic $j$
are
\[
 P=R\rho_*[E-(1-N)\End E],\qquad
 T^{\rm vir}=P-\kappa P^\vee,\qquad
 \sigma^{j+\rk P}(\det P)^{-1}.
\]
On the rotated space the virtual tangent and modified determinant are
\[
 R\rho_*[E+cE^\vee-(1-q)(1-c/q)\End E],\qquad
 c^{k+mb/2}\mathcal D(E)^{-1}\langle\det E,L/K\rangle .
\]
The counting signs are $(-1)^j$ and $(-1)^k$, respectively.
\end{lemma}
\begin{proof}
The two coordinate Koszul resolutions give the endomorphism terms,
and the section and its Serre dual give the two linear terms.
In the first presentation curve duality pairs the two halves since
$N(\kappa K/N)=\kappa K$. The chosen square root of
$\det(T^{\rm vir})^{-1}$ is $\sigma^{\rk P}/\det P$ and the ordered
external ratio contributes $\sigma^j$ after removing the counting
variable. Riemann--Roch gives $\rk P=j+m^2b$.

In the second presentation the endomorphism term with its two
Koszul factors has trivial determinant. The determinant of the
linear terms, combined with the external ratio, has underlying
square
\[
 \frac{\mathcal D(EL)}{\mathcal D(EM)\mathcal D(E)\mathcal D(E^\vee)}.
\]
For a line $H$, the identity
$\mathcal D(EH)/\mathcal D(E)=
 (\mathcal D(H)/\mathcal D(\OO))^m\langle\det E,H\rangle$
follows from the determinant definition of the Deligne pairing and
additivity (or a splitting of $E$). Curve duality gives
$\mathcal D(E^\vee)=\mathcal D(EK)$ and
$\mathcal D(M)=\mathcal D(L)$. Substitution gives the square of
$\mathcal D(E)^{-1}\langle\det E,L/K\rangle$ up to a constant line.
Taking this ordered expression chooses the root in families. Its
character exponent is $(\chi(EM)-\chi(E^\vee))/2=k+mb/2$.
\end{proof}

\begin{lemma}\label{lem:ring}
In either presentation all coefficients belong to
$q^{-D_m}K^0(B)[\sigma^{\pm1}][[q,t]]$. Multiplication by $\sigma^{mb}$
makes the powers of $\sigma$ even.
\end{lemma}
\begin{proof}
First localize only at the antidiagonal subtorus of the internal
plane. Every geometric fixed support lies over the zero section,
since the connected torus fixes each point in a finite support
fiber. The stable-pair support morphism is proper, by the projective
bundle compactification and restriction to cycles disjoint from
infinity. Thus the fixed scheme is proper over $B$, also with its
nilpotents. Each nonzero moving weight has a unit extremal Euler
coefficient, a signed line. Its Laurent expansion can be pushed
forward coefficientwise with the residual volume equivariance still
present. On $B$ this gives Laurent polynomials in $\sigma$.
Successive localization agrees with full-torus localization. It
suffices to check the bounds after inverting $1-\sigma^r$,
$r>0$, since these are non-zero-divisors even over a ring with
additive torsion.

On the original full fixed locus the box $(i,j)$ gives the
bundle $L^{-i}M^{-j}\OO(D_{ij})$ with its coordinate character.
Thus the Euler exponent is $-b\sum h_x+\sum|D_{ij}|$.
Here $\sum h_x=\sum(1+i+j)\le m(m+1)/2$ by a linear extension of the
diagram. Reflecting a negative $t$-weight in a quotient of two Euler
classes with the same bundle changes its constant by a volume unit.
There are no negative $t$ powers. On the rotated side $k=\sum |D_{ij}|$.
To bound the $q$ exponent, forget volume weights in the \emph{ranks}
of the moving tangent. Put
$V=\sum q^{j-i}$ and $P_1=\sum |D_{ij}|q^{j-i}$, and write
$X^*=X(q^{-1})$. The change from the empty tangent has rank polynomial
\[
 F=P_1-P_1^*-(1-q)(1-q^{-1})(P_1V^*-VP_1^*).
\]
Its coefficients are antisymmetric and $\sum_{i>0}i[q^i]F=P_1'(1)$.
The reflected moving Euler has this change in its lowest exponent.
The determinant root changes it by $-P_1'(1)$, while the ordered
external lines have equal degrees. The empty tangent is $-b$ times
the Hilbert tangent, with positive exponents $h_x$, giving the bound.
The two root exponents in Lemma~\ref{lem:root} differ from $-mb$ by
even integers in the $\sigma$ variable. This proves the assertion.
\end{proof}
